% !TeX root = ../../Book3.tex
% 纲要标题：## 第19章　Cohen–Macaulay 环与完全交
% 纲要位置：计划纲要.md，第 2555 行。

\chapter{Cohen–Macaulay 环与完全交}
\label{b3:ch:19}
\BookChapterNumber{b3:ch:19}

\begin{chapterabstract}
Cohen–Macaulay 性把维数与深度联系起来：参数系不仅能把支集降至有限长度，而且逐项都是非零因子。本章建立这个判据，证明局部化、商去正则序列及极大理想进完备化中的保持性质，再用 Koszul 消解说明完全交的结构。超曲面、嵌入分量和零维局部环提供了性质之间不能逆推的例子；参数商长度则使深度条件与重数计算发生直接联系。
\end{chapterabstract}

\begin{learninggoals}
\begin{itemize}
\item 区分 CM 模、极大 CM 模及非局部环上的逐点定义。
\item 证明 CM 模的参数系判据，并核对商环、局部化和完备化所需条件。
\item 从给定正则局部环与正则序列构造完全交，写出其环境环上的 Koszul 消解。
\item 用基座、参数商长度及嵌入关联素理想区分正则、完全交和 CM 性，并计算 CM 参数理想的重数。
\end{itemize}
\end{learninggoals}

设 \(k\) 为域。交叉直线的局部环 \(k[x,y]_{(x,y)}/(xy)\) 在交点不正则，但 \(x+y\) 仍是非零因子；它的一维支集因而允许长为一的正则序列。相反，在 \(k[x,y]_{(x,y)}/(x^2,xy)\) 中，极大理想杀死非零的 \(\bar x\)，同样的一维支集上连第一步也做不到。奇点是否正则，与参数能否逐项保持非零因子，是两个不同的问题。

\ProseParagraphBreak

Cohen–Macaulay 条件要求后一个问题达到维数允许的最大长度。这样一个数值条件为何能控制参数商的长度和重数，需要借助 Koszul 复形和深度来证明；完全交则提供了一批结构更强的例子。在一般环境环中，方程个数等于理想高度仍不足以保证这些方程逐项正则；正则局部环是一个重要例外，本章将把给定方程补成参数系来证明这一判据。

\ProseParagraphBreak

本章使用\chapref{b3:ch:09}的参数与维数、\chapref{b3:ch:13}的完备化、\chapref{b3:ch:15}的正则序列与 Koszul 复形，以及\chapref{b3:ch:16}的深度。正则局部环的正则参数系由\chapref{b3:ch:18}建立。Ext 与 Tor 的消解计算使用\cref{b3:lem:ext-basic-properties}和\cref{b3:lem:tor-basic-properties}；本章在完备化与基座证明中写出实际需要的复形比较。复形、投射消解及导出函子的系统讲授位置为第四卷第6—9章。

\ProseParagraphBreak
松村英之的《Commutative Ring Theory》在\cite[Theorems 17.2--17.5, pp. 134--136]{Matsumura1989CommutativeRingTheory}讨论 CM 模、参数系及完备化；正则局部环的 CM 性见\cite[Theorem 17.8, p. 137]{Matsumura1989CommutativeRingTheory}，伴随分次刻画与参数理想重数判据见\cite[Theorems 17.10--17.11, pp. 138--139]{Matsumura1989CommutativeRingTheory}。本章采用模版本陈述，并把“给定的正则序列商”与完全交局部环的内在定义区别开。

% !TeX root = ../../Book3.tex
% 纲要标题：### 19.1 Cohen–Macaulay 模
% 纲要位置：计划纲要.md，第 2557 行。

\section{Cohen–Macaulay 模}
\label{b3:sec:1901}

维数描述支集中素理想链的长度，深度则数出在极大理想中能够连续选取多少个非零因子。一般只有深度不超过维数；两者相等时，每次降维都能由一个不产生额外零化关系的方程完成。这种相等性比正则性宽松，因而能够容纳许多有奇点的环。


% 纲要内容（仅作写作计划，尚非教材正文）：
%
% * 深度等于模维数；极大 CM 模等价于 CM 模且支集维数达到环境环维数
% * 局部与非局部定义；局部 CM 模的关联素理想同维且无嵌入素理想，反向用极大理想模说明不成立
% * 参数系与正则序列的等价判据，以根理想等式和高度定理证明每个参数系正则
%

\subsection{深度等于维数}
\label{b3:subsec:1901:01}

\begin{definition}\label{b3:def:cohen-macaulay-module}
设 \((R,\mathfrak m)\) 为交换 Noether 局部环，\(M\) 为有限生成的 \(R\) 模。若 \(M\ne0\) 且
\[
\operatorname{depth}_R M=\dim_R M,
\]
则称 \(M\) 为 \term[CohenMacaulaymo]{Cohen–Macaulay 模}[Cohen--Macaulay module]，简称 CM 模。也约定零模为 CM 模，但不把零模代入上述有限数值等式。若 \(R\) 作为自身的模为 CM 模，则称 \(R\) 为 \term[CohenMacaulayhuan]{Cohen–Macaulay 环}[Cohen--Macaulay ring]。对于非零的 \(M\)，若 \(\operatorname{depth}_RM=\dim R\)，则称 \(M\) 为\term[jidaCohenMacaulaymo]{极大 Cohen–Macaulay 模}[maximal Cohen--Macaulay module]。
\end{definition}

这里 \(\dim_RM=\dim(R/\operatorname{Ann}_R M)\)，它未必等于 \(\dim R\)。由\cref{b3:cor:depth-dimension-bound}，
\[
\operatorname{depth}_RM\le\dim_RM\le\dim R,
\]
所以极大 CM 模恰好是支集维数达到 \(\dim R\) 的 CM 模。“极大”指深度达到环境环的维数，并非按子模包含关系取极大者。例如，对域 \(k\)，取 \(R=k[x,y]_{(x,y)}\) 和 \(M=R/(x)\)，则
\[
\operatorname{depth}_RM=\dim_RM=1<2=\dim R.
\]
因此 \(M\) 为 CM 模，却不是极大 CM 模。

\begin{proposition}\label{b3:prop:cm-associated-dimensions}
设 \((R,\mathfrak m)\) 为交换 Noether 局部环，\(M\ne0\) 为\(R\) 上有限生成的 CM 模，\(d=\dim_RM\)。则每个 \(\mathfrak p\in\operatorname{Ass}_RM\) 都满足 \(\dim(R/\mathfrak p)=d\)，并且
\(\operatorname{Ass}_RM=\operatorname{Min}\operatorname{Supp}_RM\)。
反向结论不成立。具体地，若 \(k\) 为域、\(R_0=k[x,y]_{(x,y)}\)、\(\mathfrak m_0=(x,y)R_0\)，则
\[
\operatorname{Ass}_{R_0}\mathfrak m_0=\{(0)\},\qquad
\dim_{R_0}\mathfrak m_0=2,\qquad
\operatorname{depth}_{R_0}\mathfrak m_0=1.
\]
因此 \(\mathfrak m_0\) 不是 CM 模；但对每个非极大素理想 \(\mathfrak p\subset R_0\)，均有 \((\mathfrak m_0)_{\mathfrak p}=(R_0)_{\mathfrak p}\)，它是 CM 模。
\end{proposition}
\begin{proof}
对 \(\mathfrak p\in\operatorname{Ass}_RM\)，\cref{b3:cor:associated-depth-bound}与 \(\mathfrak p\in\operatorname{Supp}_RM\) 给出
\[
d=\operatorname{depth}_RM\le\dim(R/\mathfrak p)\le\dim_RM=d.
\]
因此每个关联素理想的商环都达到支集的最大维数。若 \(\mathfrak p\) 严格包含支集中的极小素理想 \(\mathfrak q\)，则在从 \(\mathfrak p\) 到 \(\mathfrak m\) 的一条最长链前添上 \(\mathfrak q\)，得到
\(\dim(R/\mathfrak q)\ge\dim(R/\mathfrak p)+1>d\)，矛盾。反向包含由\cref{b3:prop:minimal-primes-associated}成立。

对反例，因为 \(R_0\) 是整环，\(\mathfrak m_0\) 中每个非零元素的零化子都是零，故 \(\operatorname{Ass}_{R_0}\mathfrak m_0=\{(0)\}\)，且 \(\operatorname{Ann}_{R_0}\mathfrak m_0=0\) 给出 \(\dim_{R_0}\mathfrak m_0=\dim R_0=2\)。乘 \(x\) 在 \(\mathfrak m_0\) 上单射；但在 \(\mathfrak m_0/x\mathfrak m_0\) 中，\(x\) 的类非零，否则约去 \(x\) 便得 \(1\in\mathfrak m_0\)。这个类被 \(x,y\) 同时零化，所以\cref{b3:cor:depth-zero-socle}给出 \(\operatorname{depth}_{R_0}(\mathfrak m_0/x\mathfrak m_0)=0\)。再由\cref{b3:cor:regular-quotient-depth}，\(\operatorname{depth}_{R_0}\mathfrak m_0=1\)，所以 \(\mathfrak m_0\) 不是 CM 模。

若 \(\mathfrak p\ne\mathfrak m_0\)，则 \(x,y\) 中至少一个不属于 \(\mathfrak p\)，因而 \((\mathfrak m_0)_{\mathfrak p}=(R_0)_{\mathfrak p}\)。\cref{b3:thm:regular-local-localization}使 \((R_0)_{\mathfrak p}\) 仍为正则局部环；由\cref{b3:thm:regular-local-characterizations}，它的深度等于维数，故为 CM 环。
\end{proof}

所以局部 CM 模的各不可约分量维数相同，而且没有嵌入关联素理想：较低维的嵌入关联素理想与维数不同的极小分量都被排除了。反过来，仅知道没有嵌入关联素理想并不足以得到 CM 性；在维数至少二时，还要检查商去一个非零因子后是否仍能继续。

\subsection{Cohen–Macaulay 性的局部与非局部定义}
\label{b3:subsec:1901:02}

\begin{definition}\label{b3:def:global-cohen-macaulay}
设 \(R\) 为交换 Noether 环，\(M\) 为有限生成的 \(R\) 模。若对每个 \(\mathfrak p\in\operatorname{Supp}_RM\)，局部环 \(R_{\mathfrak p}\) 上的模 \(M_{\mathfrak p}\) 都为 CM 模，则称 \(M\) 为 CM 模。取 \(M=R\) 得到非局部环的 CM 性定义。
\end{definition}

局部定义只比较一个极大理想处的两个整数，非局部定义则逐点比较。若局部 CM 性在进一步局部化后保持，检查所有极大理想便足够；这一保持性质见\cref{b3:thm:cm-localization}。

\ProseParagraphBreak
对有限直积 \(R_1\times\cdots\times R_s\)，每个素理想只来自一个因子，故该环为 CM 环当且仅当每个因子为 CM 环。不同因子的维数可以不同。例如，对域 \(k\)，\(k\times k[t]\) 的两个连通分支分别为 \(\operatorname{Spec}k\) 和 \(\operatorname{Spec}k[t]\)，维数为零和一；它们在每个素点处仍是 CM 的局部环。因此，局部 CM 模的各不可约分量等维，并不要求一个非局部 CM 环的不同连通分支具有相同维数。

\subsection{参数系与正则序列}
\label{b3:subsec:1901:03}

\begin{definition}\label{b3:def:module-system-parameters}
设 \((R,\mathfrak m)\) 为交换 Noether 局部环，\(M\ne0\) 为有限生成的 \(R\) 模，\(d=\dim_RM\)。若 \(x_1,\ldots,x_d\in\mathfrak m\) 且
\(M/(x_1,\ldots,x_d)M\) 为有限长度模，则称此序列为 \(M\) 的\term[mochanshuxi]{参数系}[system of parameters for a module]。对 \(d=0\)，采用空序列。
\end{definition}

令 \(A=R/\operatorname{Ann}_RM\)。对任意理想 \(I\subseteq R\) 及素理想 \(\mathfrak p\)，局部化给出
\[
(M/IM)_{\mathfrak p}
=M_{\mathfrak p}/I_{\mathfrak p}M_{\mathfrak p}.
\]
若 \(I\not\subseteq\mathfrak p\)，则 \(I_{\mathfrak p}=R_{\mathfrak p}\)，这个商为零。若 \(I\subseteq\mathfrak p\) 且 \(M_{\mathfrak p}\ne0\)，则 \(I_{\mathfrak p}\subseteq\mathfrak pR_{\mathfrak p}\)，Nakayama 引理保证这个有限生成商模非零。因此
\[
\operatorname{Supp}(M/IM)=\operatorname{Supp}(M)\cap V(I)
\]
成立。一个非零有限生成局部模若有有限长度，其组成因子均为剩余域，所以支集只有闭点。反过来，若支集只有闭点，则模的零化理想所给商环为零维 Noether 局部环，由\cref{b3:thm:noether-dimension-zero,b3:thm:artin-finite-length}，该商环及其有限生成模均有有限长度。因此参数系等价于其像在 \(A\) 中生成一个极大理想准素理想；\cref{b3:thm:system-of-parameters}保证它存在，并保证少于 \(d\) 个元素不能达到这一要求。

\begin{theorem}\label{b3:thm:cm-parameters}
设 \((R,\mathfrak m)\) 为交换 Noether 局部环，\(M\ne0\) 为有限生成的 \(R\) 模，\(d=\dim_RM\)。下列条件等价：
\begin{equivalences}
\item \(M\) 为 CM 模；
\item \(M\) 的每个参数系均为 \(M\) 正则序列；
\item \(M\) 的某个参数系为 \(M\) 正则序列。
\end{equivalences}
此外，在 CM 情形中，任意参数系的前 \(r\) 项所给商模仍为 CM 模，维数为 \(d-r\)。
\end{theorem}
\begin{proof}
先证明 CM 模的每个参数系正则。关键是排除第一个参数为零因子：若它在一个 \(d\) 维关联分量上已经为零，余下的 \(d-1\) 个参数就必须单独把这一分量切到闭点，而高度定理不允许 \(d-1\) 个方程完成这样的降维。

对 \(d\) 归纳；\(d=0\) 时没有元素需要检验。取参数系 \(x_1,\ldots,x_d\)。若 \(x_1\) 是 \(M\) 的零因子，则它属于某个 \(\mathfrak p\in\operatorname{Ass}_RM\)，而\cref{b3:prop:cm-associated-dimensions}给出 \(\dim R/\mathfrak p=d\)。参数系满足
\[
\sqrt{(x_1,\ldots,x_d)+\operatorname{Ann}_RM}=\mathfrak m.
\]
由于 \(\operatorname{Ann}_RM\subseteq\mathfrak p\subseteq\mathfrak m\)，把等式中的理想加上 \(\mathfrak p\) 再取商，并用 \(x_1\in\mathfrak p\)，得到
\[
\sqrt{(x_2,\ldots,x_d)(R/\mathfrak p)}
=\mathfrak m/\mathfrak p.
\]
因此在 Noether 局部环 \(R/\mathfrak p\) 中，极大理想在一个由 \(d-1\) 个元素生成的理想之上极小。\cref{b3:thm:height-theorem}给出
\[
d=\dim(R/\mathfrak p)
=\operatorname{ht}_{R/\mathfrak p}(\mathfrak m/\mathfrak p)
\le d-1,
\]
矛盾。

所以 \(x_1\) 避开所有关联素理想，是 \(M\) 上的非零因子。\cref{b3:cor:regular-quotient-depth}和\cref{b3:prop:regular-quotient-dimension}分别给出
\[
\operatorname{depth}(M/x_1M)=d-1=\dim(M/x_1M).
\]
余下的元素为这个非零商模的参数系，归纳即可；逐次应用同一论证也得到附加结论。

每个参数系正则当然推出存在一个正则参数系。若存在长度为 \(d\) 的正则参数系，则深度至少为 \(d\)，而深度不超过维数，故两者相等。
\end{proof}

这个判据的主要内容在于“每个”：CM 性使任意参数系都逐项正则。在一般局部环中，参数系只保证把支集降至闭点；在 CM 模中，它同时保证每次乘法都单射，任意一组参数都能连续完成维数所允许的切割。

\ProseParagraphBreak
这并不要求环约化。例如，对域 \(k\)，令 \(B=k[x,y]_{(x,y)}/(x^2)\)，则 \(\dim B=1\)，\(\bar y\) 是参数又是非零因子，所以 \(\operatorname{depth}B=1\)。但 \(\bar x\ne0\) 且 \(\bar x^2=0\)。CM 性控制的是正则序列的长度，不能用来排除幂零元，也不能据此判断不可约性或正则性。

% !TeX root = ../../Book3.tex
% 纲要标题：### 19.2 Cohen–Macaulay 性的保持与正则环
% 纲要位置：计划纲要.md，第 2563 行。

\section{Cohen–Macaulay 性的保持与正则环}
\label{b3:sec:1902}

要把 Cohen–Macaulay 性用于计算，需要知道它在商、局部化与完备化下怎样变化。正则序列同时降低深度与维数，另外两种操作则通过局部不等式及忠实平坦性保持二者相等。


% 纲要内容（仅作写作计划，尚非教材正文）：
%
% * 商去正则序列同时降低深度与模维数；商环与原环下的支集自然对应
% * 用支集中素链拼接和深度不等式证明 CM 性局部化及支集维数等式；区分所有极大理想处检验与稠密开集检验
% * 用各项有限秩自由消解、Hom 基变换及平坦张量与上同调交换证明极大理想进完备化保持深度，忠实性反映零性
% * 正则局部环为 CM 环；非局部 CM 性只须检验所有极大理想
%

\subsection{商去正则元素}
\label{b3:subsec:1902:01}

\begin{theorem}\label{b3:thm:cm-regular-quotient}
设 \((R,\mathfrak m)\) 为交换 Noether 局部环，\(M\ne0\) 为有限生成的 \(R\) 模，\(r\ge0\) 为整数，\(x_1,\ldots,x_r\in\mathfrak m\) 为 \(M\) 正则序列。令 \(I=(x_1,\ldots,x_r)\)、\(\overline M=M/IM\)。则
\[
M\;\text{为 CM 模}\quad\Longleftrightarrow\quad
\overline M\;\text{为 CM 模},
\qquad \dim\overline M=\dim M-r.
\]
右边既可把 \(\overline M\) 看成 \(R\) 模，也可看成 \(R/(x_1,\ldots,x_r)\) 模。更一般地，对任意非零有限生成的 \(R\) 模 \(L\) 及理想 \(J\subseteq\operatorname{Ann}_RL\)，\(L\) 作为 \(R\) 模与作为 \(R/J\) 模的深度、维数及 CM 性均相同。
\end{theorem}
\begin{proof}
逐次商去一个正则元素时，深度与维数同步下降，所以只需比较它们的差，并确认每一步的商仍非零。设当前逐次商为 \(N\ne0\)，下一元素为 \(x\in\mathfrak m\)。若 \(N/xN=0\)，则 \(N=xN\)，Nakayama 引理给出 \(N=0\)，矛盾。因此各个逐次商非零。对每一步，\cref{b3:cor:regular-quotient-depth}使深度恰减一，\cref{b3:prop:regular-quotient-dimension}使维数恰减一，故
\[
\dim\overline M-\operatorname{depth}_R\overline M
=\dim M-\operatorname{depth}_RM.
\]
这个差为零恰好表示 CM 性，因而商去正则序列既不会消除原有的差，也不会产生新的差。

最后比较零化理想取商前后的底环。对 \(J\subseteq\operatorname{Ann}_RL\)，\(R\) 的作用通过 \(R/J\) 分解，两种作用给出相同的子模。素理想对应
\[
\operatorname{Spec}(R/J)\xrightarrow{\sim}V(J)\subseteq\operatorname{Spec}R,
\qquad \mathfrak p/J\longmapsto\mathfrak p
\]
是保持包含关系的自然同胚，并把 \(L\) 在两种底环下的支集对应起来，所以模维数相同。来自 \(\mathfrak m\) 的元素和它在 \(\mathfrak m/J\) 中的像作用为同一个乘法算子；逐次商也不改变。因此正则序列在取像和提升时保持正则性，两种底环下的深度也相同。取 \(L=\overline M\)、\(J=I\)，就得到正则商在两种底环下的相同结论。
\end{proof}

非零因子假设不能删除。设 \(k\) 为域，取 \(R=k[x,y]_{(x,y)}/(x^2,xy)\)，则商去 \(x\) 后得到正则的一维局部环 \(k[y]_{(y)}\)，而原环的非零元 \(x\) 被整个极大理想零化，深度为零。这时商环的 CM 性不能反映原环的 CM 性。

\subsection{Cohen–Macaulay 性的局部化与完备化}
\label{b3:subsec:1902:02}

\begin{theorem}\label{b3:thm:cm-localization}
设 \((R,\mathfrak m)\) 为交换 Noether 局部环，\(M\ne0\) 为有限生成的 \(R\) 模且为 CM 模。对每个 \(\mathfrak p\in\operatorname{Supp}_RM\)，\(M_{\mathfrak p}\) 都是 \(R_{\mathfrak p}\) 上的 CM 模，并且
\[
\dim_RM=\dim_{R_{\mathfrak p}}M_{\mathfrak p}+\dim(R/\mathfrak p).
\]
更一般地，设 \(A\) 为交换 Noether 环，\(N\) 为有限生成的 \(A\) 模，\(S\subseteq A\) 为乘法集；若 \(N_{\mathfrak n}\) 对每个极大理想 \(\mathfrak n\subseteq A\) 都是 CM 模，则 \(N\) 为 CM 模，且 \(S^{-1}N\) 在 \(S^{-1}A\) 上仍为 CM 模。
\end{theorem}
\begin{proof}
素链拼接给出局部维数与 \(\dim R/\mathfrak p\) 之和的上界，深度局部化不等式给出相应的下界；CM 假设使两边相遇，于是同时得到局部化后的 CM 性和维数等式。

有限生成性给出 \(\operatorname{Supp}_RM=V(\operatorname{Ann}_RM)\)，因此每个包含 \(\mathfrak p\) 的素理想仍在支集中。在 \(\operatorname{Supp}_{R_{\mathfrak p}}M_{\mathfrak p}\) 中取一条长度为 \(\dim M_{\mathfrak p}\)、以 \(\mathfrak pR_{\mathfrak p}\) 为末项的链，收缩到 \(R\) 后得到支集中止于 \(\mathfrak p\) 的链。又因 \(R/\mathfrak p\) 是局部整环，可取一条长度为 \(\dim R/\mathfrak p\)、从零理想到 \(\mathfrak m/\mathfrak p\) 的链；提升后得到从 \(\mathfrak p\) 到 \(\mathfrak m\) 的链，其各项也在支集中。把这两条链在 \(\mathfrak p\) 处拼接，便得
\[
\dim M_{\mathfrak p}+\dim R/\mathfrak p\le\dim M.
\]
另一方面，\cref{b3:thm:depth-localization}给出
\[
\dim M=\operatorname{depth}_RM
\le\operatorname{depth}_{R_{\mathfrak p}}M_{\mathfrak p}+\dim R/\mathfrak p
\le\dim M_{\mathfrak p}+\dim R/\mathfrak p
\le\dim M.
\]
所以所有不等式都是等式，尤其局部化后的深度等于维数，并且所述维数公式成立。一般情形中，对 \(\mathfrak p\in\operatorname{Supp}_AN\)，选极大理想 \(\mathfrak n\supseteq\mathfrak p\)，则 \(N_{\mathfrak n}\ne0\) 为 CM 模。在 \(A_{\mathfrak n}\) 中进一步局部化到 \(\mathfrak pA_{\mathfrak n}\)，所得环和模分别为 \(A_{\mathfrak p}\) 和 \(N_{\mathfrak p}\)，已证结论因而适用。局部化环的素点又与原环中避开乘法集的素点对应，故 \(S^{-1}N\) 在其支集各点也为 CM 模。
\end{proof}

检验所有极大理想处的局部化足够，检验一个稠密开集却可能遗漏 CM 性失败的点。\cref{b3:prop:cm-associated-dimensions}中的 \(R_0=k[x,y]_{(x,y)}\)、\(\mathfrak m_0=(x,y)R_0\) 已给出这样的反例：\(\mathfrak m_0\) 的深度为一、维数为二，但每个非极大素点处的局部化都是 CM 模。开集
\[
\operatorname{Spec}R_0\setminus\{\mathfrak m_0\}=D(x)\cup D(y)
\]
包含整环 \(R_0\) 的一般点，因而稠密；\(\mathfrak m_0\) 在这里处处为 CM 模，失败仍留在被删去的闭点。

\begin{theorem}\label{b3:thm:cm-completion}
设 \((R,\mathfrak m)\) 为交换 Noether 局部环，\(M\ne0\) 为有限生成的 \(R\) 模。记 \(\widehat R\) 和 \(\widehat M\) 为极大理想进完备化，则
\[
\operatorname{depth}_{\widehat R}\widehat M=\operatorname{depth}_RM,
\qquad M\;\text{为 CM 模}\Longleftrightarrow\widehat M\;\text{为 CM 模}.
\]
\end{theorem}
\begin{proof}
由深度的 Ext 刻画，目标是比较完备化前后 Ext 的首次非零次数。为此，将剩余域的自由消解张量到完备环，利用有限秩逐项比较 Hom，再用平坦性比较上同调，最后用忠实性检测非零项。

令 \(k=R/\mathfrak m\)。由\cref{b3:thm:completion-flat,b3:prop:completed-local-ring}，\(\widehat R\) 是 Noether 局部环，极大理想为 \(\mathfrak m\widehat R\)，且 \(R\to\widehat R\) 忠实平坦，剩余域仍为 \(k\)。\cref{b3:thm:completion-tensor}给出 \(\widehat M\cong M\otimes_R\widehat R\)，故 \(\widehat M\) 是非零有限生成的 \(\widehat R\) 模。

取 \(k\) 的自由消解 \(F_\bullet\to k\)，其中各项均为有限秩自由模。构造时先用有限秩自由模满射到有限生成的 \(k\)，其核由 Noether 性仍有限生成，再重复这一过程。消解可能无限延伸，但每项秩有限。\(\widehat R\) 的平坦性使张量后的复形仍消解 \(k\otimes_R\widehat R\cong k\)。有限秩保证逐项有与微分相容的自然同构
\[
\operatorname{Hom}_R(F_i,M)\otimes_R\widehat R
\cong\operatorname{Hom}_{\widehat R}(F_i\otimes_R\widehat R,\widehat M).
\]
具体地，若 \(F_i\cong R^{a_i}\)，两边均为 \(\widehat M^{a_i}\)，这里只需张量积与有限直和交换。若换成无限秩自由模，Hom 会产生无限乘积，而张量积一般不与无限乘积交换；Noether 性在这里保证了所需的有限性。

令 \(C^\bullet=\operatorname{Hom}_R(F_\bullet,M)\)。平坦性还说明 \(-\otimes_R\widehat R\) 是正合函子，它保持核、像及其商，从而对每个 \(i\ge0\) 有自然同构
\[
H^i(C^\bullet)\otimes_R\widehat R
\cong H^i(C^\bullet\otimes_R\widehat R).
\]
结合前面的逐项同构，得到
\[
\operatorname{Ext}_R^i(k,M)\otimes_R\widehat R
\cong\operatorname{Ext}_{\widehat R}^i(k,\widehat M).
\]
前面传递消解的正合性和比较上同调只使用平坦性。忠实性则保证，对任意 \(R\) 模 \(E\)，
\[
E\otimes_R\widehat R=0\quad\Longleftrightarrow\quad E=0.
\]
将其用于各个 \(\operatorname{Ext}_R^i(k,M)\)，可知两环下 Ext 的零项和非零项出现于相同次数。由\cref{b3:thm:depth-ext}，首次非零次数相同便给出深度相同。\cref{b3:thm:completion-dimension}已证明有限生成模的维数也相同，遂得 CM 性的等价。
\end{proof}

这里使用的完备化是 Noether 局部环的极大理想进完备化。若在任意非局部环中选一个理想完备化，映射未必忠实；它可能完全丢掉某些素点，因而不能不加条件地用完备环检测原环的全部 CM 性。

\subsection{正则局部环的 Cohen–Macaulay 性}
\label{b3:subsec:1902:03}

\begin{corollary}\label{b3:cor:regular-ring-cm}
每个交换 Noether 正则局部环都是 CM 环。若交换 Noether 环的所有素点局部环均正则，则该环为 CM 环。
\end{corollary}
\begin{proof}
\cref{b3:thm:regular-local-characterizations}说明：一个 \(d\) 维正则局部环的极大理想可由长度为 \(d\) 的正则序列生成。因此深度至少为 \(d\)，结合深度不超过维数即得相等。非局部结论按定义逐点应用。
\end{proof}

设 \(k\) 为域，\(d\ge1\) 为整数。在 \(k[\![x_1,\ldots,x_d]\!]\) 及 \(k[x_1,\ldots,x_d]_{(x_1,\ldots,x_d)}\) 中，变量依次正则且最后的商为 \(k\)。把这些环商去一个非零非单位，就得到维数减一的 CM 环，但商环可能有奇点，甚至不约化。下一节将统一处理由多个正则元素给出的商。

% !TeX root = ../../Book3.tex
% 纲要标题：### 19.3 完全交
% 纲要位置：计划纲要.md，第 2569 行。

\section{完全交}
\label{b3:sec:1903}

在一般环境环中，仅数方程和余维，不能保证每个方程在前面的商上仍为非零因子。在正则局部环中则有更强的判据：真理想的生成元数若等于其高度，这组生成元就是正则序列。完全交的内在定义要求，环的完备化能写成完备正则局部环的正则序列商。这样的方程组有 Koszul 消解，因而维数与深度可以同时计算。


% 纲要内容（仅作写作计划，尚非教材正文）：
%
% * 完全交局部环的完备化定义；用极大理想极小生成组的第一 Koszul 同调比较任意正则局部环境，证明完全交等价于当前核的最小生成元数等于高度，并区分极小与冗余方程组；实际正则环境中的局部化直接证明，一般完备化定义下的局部化注明所用平坦局部下降定理
% * Koszul 消解与 Tor 的变量交换；CM 模在正则局部环境中的投射维数等于支集余维；区分环境环有限消解与商环周期消解
% * 用两套 Koszul 消解计算同一 Tor，证明零维完全交的基座一维；平方零族给出 CM 而非完全交的反例
% * 正则蕴含完全交，完全交蕴含 CM；集中展示逆向反例及 CM 与所有 Serre 条件的等价关系
%

\subsection{正则环中正则序列定义的商}
\label{b3:subsec:1903:01}

设 \((S,\mathfrak n)\) 为正则局部环。给定 \(c\) 个方程后，若它们生成的真理想恰有高度 \(c\)，便能把这组方程补成 \(S\) 的完整参数系。\(S\) 是 CM 环，因而整组参数正则，原来的方程也就逐项正则。下面的\cref{b3:prop:given-complete-intersection}把这个判据与正则序列商的性质一起陈述。这里理想的高度仍取包含它的素理想高度的最小值；证明先计算商环的维数，以确定还需要补几个参数。

\begin{definition}\label{b3:def:local-complete-intersection}
设 \((R,\mathfrak m)\) 为交换 Noether 局部环，\(\widehat R\) 为其 \(\mathfrak m\) 进完备化。若存在完备正则局部环 \((S,\mathfrak n)\) 及 \(S\) 正则序列 \(f_1,\ldots,f_c\in\mathfrak n\)，使
\[
\widehat R\cong S/(f_1,\ldots,f_c),
\]
则称 \(R\) 为\term[wanquanjiaojubuhuan]{完全交局部环}[complete intersection local ring]。允许 \(c=0\)。对交换 Noether 环 \(A\)，若每个素理想 \(\mathfrak p\subseteq A\) 处的局部环 \(A_{\mathfrak p}\) 都是完全交局部环，则称 \(A\) 为\term[jubuwanquanjiaohuan]{局部完全交环}[locally complete intersection ring]。
\end{definition}

完全交性要求存在某个正则序列展示。在一个展示 \(R=S/I\) 中，\(I\) 的某组可能冗余的生成元不正则，不能仅据此否定完全交性。例如，重复写出同一个方程便会破坏正则性，却没有改变商环。不过，当 \(S\) 为正则局部环，且所取的是 \(I\) 的极小生成组时，情况不同：\cref{b3:thm:ci-presentation-independence}将证明，\(R\) 完全交当且仅当这组生成元正则；换用其他正则局部环境不能消除这一障碍。定义从 \(R\) 自身的完备化出发，因此不要求预先给出正则局部环到 \(R\) 的满射；完备局部环的环境构造与 Cohen 结构定理有关，见本卷附录。\cref{b3:prop:ci-localization}还将说明完全交局部环在其他素点处仍完全交，从而与一般环的局部完全交定义相容。

\ProseParagraphBreak
任意展示中的序列长度 \(c\) 也不能立刻看作 \(R\) 的内在数值。例如，对任意域 \(k\)，有
\[
k[\![x]\!]/(x^2)\cong k[\![x,T]\!]/(T,x^2).
\]
左边的正则序列长一，右边的正则序列长二，商环却相同；增加变量并增加消去该变量的关系，会同时增加环境维数与关系数。本章的局部完全交性是环自身在各素点的性质。代数几何中的局部完全交态射或正则嵌入则还涉及一个指定态射或嵌入，使用这些术语时须说明所指对象。

\begin{proposition}\label{b3:prop:given-complete-intersection}
设 \((S,\mathfrak n)\) 为交换 Noether 正则局部环，\(I=(f_1,\ldots,f_c)\subsetneq S\)，\(R=S/I\)。若 \(f_1,\ldots,f_c\) 为 \(S\) 正则序列，则 \(R\) 为完全交局部环，且 \(\dim R=\dim S-c\)。若 \(\operatorname{ht}I=c\)，则给定的 \(f_1,\ldots,f_c\) 必为 \(S\) 正则序列，因而 \(R\) 同样是维数为 \(\dim S-c\) 的完全交局部环。允许 \(c=0\)，此时 \(I=0\)。
\end{proposition}
\begin{proof}
先证明高度判据。记 \(d=\dim S\)。\cref{b3:thm:regular-local-characterizations}给出 \(\operatorname{depth}S=d\)，而\cref{b3:thm:depth-localization}与素链拼接给出，对每个 \(\mathfrak p\in\operatorname{Spec}S\)，
\[
d\le\operatorname{depth}S_{\mathfrak p}+\dim S/\mathfrak p
\le\operatorname{ht}\mathfrak p+\dim S/\mathfrak p\le d.
\]
中间一步使用 \(\operatorname{depth}S_{\mathfrak p}\le\dim S_{\mathfrak p}=\operatorname{ht}\mathfrak p\)。全部不等式于是成为等号，得到
\[
\operatorname{ht}\mathfrak p+\dim S/\mathfrak p=d.
\]
对包含 \(I\) 的素理想取商环维数的最大值，也就是对高度取最小值，便得
\(\dim S/I=d-\operatorname{ht}I=d-c\)。在 \(S/I\) 中选取 \(d-c\) 个参数 \(\bar g_1,\ldots,\bar g_{d-c}\)，并提升至 \(\mathfrak n\)。参数系的根条件给出
\[
\sqrt{(f_1,\ldots,f_c,g_1,\ldots,g_{d-c})}=\mathfrak n.
\]
因此这 \(d\) 个元素构成 \(S\) 的参数系。\(S\) 为 CM 环，\cref{b3:thm:cm-parameters}使整组参数正则，其前 \(c\) 项恰为给定方程。这也包括空序列的情形。

现在假定给定序列正则，不论它来自原假设还是刚证的高度判据。维数等式由\cref{b3:prop:regular-quotient-dimension}逐次应用得到。\cref{b3:thm:completion-exact}和\cref{b3:prop:completion-quotients}给出
\(\widehat R\cong\widehat S/(f_1,\ldots,f_c)\widehat S\)。\(S\to\widehat S\) 的平坦性保持各逐次商上的单射乘法；忠实平坦性则保证这些非零的商张量后仍非零。因此原序列在 \(\widehat S\) 上仍正则。再由\cref{b3:thm:regular-local-completion}，\(\widehat S\) 为完备正则局部环，所以上述展示符合完全交的定义。
\end{proof}

这个判据的关键是 CM 参数系定理。一般环境环缺少这项性质，不能只凭方程数与高度相等推出正则性。重复写一个方程也不会增加理想的高度，因而不能满足正则局部环中的这一判据。

\ProseParagraphBreak
设 \(k\) 为域。例如 \(k[\![x,y,z]\!]/(xy-z^2)\) 是二维完全交局部环：\(k[\![x,y,z]\!]\) 是三维正则局部整环，\(xy-z^2\) 是一个非零非单位。相比之下，\(k[\![x,y]\!]/(x^2,xy)\) 的两个给定方程不构成正则序列；在商去 \(x^2\) 后，非零的 \(x\) 被 \(xy\) 零化。

\subsection{完全交的 Koszul 消解}
\label{b3:subsec:1903:02}

\begin{proposition}\label{b3:prop:ci-koszul-resolution}\label{b3:prop:cm-pd-codimension}
设 \((S,\mathfrak n)\) 为交换 Noether 正则局部环，\(n=\dim S\)，\(k=S/\mathfrak n\) 为其剩余域。对任意非零有限生成的 \(S\) 模 \(M\)，令 \(d=\dim_SM\)，则
\[
M\;\text{为 CM 模}\quad\Longleftrightarrow\quad
\operatorname{pd}_SM=n-d.
\]
特别地，若 \(f_1,\ldots,f_c\in\mathfrak n\) 为 \(S\) 正则序列，\(R=S/(f_1,\ldots,f_c)\)，则 Koszul 复形 \(K(f_1,\ldots,f_c;S)\) 为 \(R\) 的极小自由 \(S\) 消解，其第 \(i\) 项的秩为 \(\binom ci\)。因而
\[
\dim_k\operatorname{Tor}_i^S(k,R)=\binom ci\quad(0\le i\le c),
\qquad\operatorname{pd}_S R=c.
\]
\end{proposition}
\begin{proof}
\cref{b3:thm:regular-local-homological-criterion}保证 \(M\) 的投射维数有限，\cref{b3:thm:regular-local-characterizations}给出 \(\operatorname{depth}S=n\)。因此可以应用\cref{b3:thm:auslander-buchsbaum}，得到
\[
\operatorname{pd}_SM=n-\operatorname{depth}_SM.
\]
右端等于 \(n-d\)，恰好意味着 \(\operatorname{depth}_SM=d\)，即 \(M\) 为 CM 模。

对所给的正则序列，\cref{b3:thm:koszul-regular-resolution}把其 Koszul 复形识别为 \(R\) 的自由消解。该复形第 \(i\) 项为 \(\bigwedge^iS^c\)。若 \(e_1,\ldots,e_c\) 为 \(S^c\) 的标准基，则这一项的基为
\[
e_{j_1}\wedge\cdots\wedge e_{j_i}
\qquad(1\le j_1<\cdots<j_i\le c),
\]
即从 \(c\) 个方程中选取 \(i\) 个所得到的 \(i\) 元子集，故其秩为 \(\binom ci\)。微分矩阵的各项是 \(0\) 或 \(\pm f_j\)，均属于 \(\mathfrak n\)，故此消解极小。与 \(k\) 张量后所有微分为零，首先得到
\[
\operatorname{Tor}_i^S(R,k)\cong k^{\binom ci}.
\]
再由\cref{b3:lem:tor-basic-properties}的自然同构
\(\operatorname{Tor}_i^S(R,k)\cong\operatorname{Tor}_i^S(k,R)\)，得到所列 Tor 维数。最高次 Tor 在 \(c\) 次非零，而自由消解长度为 \(c\)，故投射维数恰为 \(c\)。
\end{proof}

第一项等价把 CM 性翻译为环境环上的消解长度：若模的支集在正则局部环境中有余维 \(n-d\)，那么 CM 性恰好要求其投射维数达到这个值。正则序列商的 Koszul 消解实现了长度为 \(c\) 的这种情形。不过，CM 商环也可以有不是 Koszul 形状的极小消解，长度等于余维并不保证关系由正则序列生成。

\ProseParagraphBreak
以上的 \(\operatorname{pd}_SR=c\) 和 \(\operatorname{Tor}_i^S(k,R)\) 都以 \(S\) 为底环，所消解的是 \(S\) 模 \(R\)。改以 \(R\) 为底环时，\(R\) 本身自由，而剩余域 \(k\) 的投射维数可能无穷。例如，对域 \(k\) 和整数 \(n\ge2\)，后面的\cref{b3:prop:ci-cm-counterexample-families}给出 \(k[\![x]\!]/(x^n)\) 上的这种周期消解。完全交性提供正则环境环上的有限消解；商环剩余域消解的有限性则由\cref{b3:thm:regular-local-homological-criterion}刻画，对应于商环本身的正则性。

\subsection{完全交、Cohen–Macaulay 性与正则性的关系}
\label{b3:subsec:1903:03}

\begin{theorem}\label{b3:thm:complete-intersection-cm}
对于交换 Noether 局部环，有
\[
\text{正则}\ \overset{c=0}{\Longrightarrow}\ \text{完全交}
\ \overset{\text{正则序列商}}{\Longrightarrow}\ \text{Cohen--Macaulay}.
\]
\end{theorem}
\begin{proof}
正则环的完备化仍正则，取空序列即可得到第一项蕴含。若 \(R\) 为完全交局部环，则按定义 \(\widehat R=S/(f_1,\ldots,f_c)\)。\(S\) 为 CM 环，\cref{b3:thm:cm-regular-quotient}使该商为 CM 环，再由\cref{b3:thm:cm-completion}下降至 \(R\)。
\end{proof}

\(k[\![x]\!]/(x^2)\) 是完全交而非正则：其维数为零，极大理想却非零。局部环 \(B=k[\![x,y]\!]/(x^2,xy)\) 则不是 CM 环，因而不是完全交：\(B\) 的维数为一，而非零的 \(\bar x\) 被极大理想 \((\bar x,\bar y)\) 零化，所以基座非零，由\cref{b3:cor:depth-zero-socle}得到 \(\operatorname{depth}B=0\)。这里使用的是环自身的 CM 障碍，而非仅仅检查当前方程组。

\ProseParagraphBreak
要给出 CM 而非完全交的例子，可以在零维局部环上使用基座。设 \(k\) 为域，先看 \(C=k[\![x,y]\!]/(x^2,y^2)\)。它有基 \(1,x,y,xy\)，且对 \(a=a_0+a_1x+a_2y+a_3xy\)，条件 \(xa=ya=0\) 恰好要求 \(a_0=a_1=a_2=0\)。因此 \(\operatorname{Soc}(C)=k\cdot xy\) 为一维。相比之下，\(k[x,y]/(x,y)^2\) 的极大理想平方为零，其基座为二维的 \(k\cdot x\oplus k\cdot y\)。下面说明，前一种一维现象是所有零维完全交局部环都必须满足的条件。

\begin{lemma}\label{b3:lem:artinian-ci-socle}
设 \((R,\mathfrak m)\) 为零维交换 Noether 完全交局部环，\(k=R/\mathfrak m\) 为其剩余域。则其基座
\(\operatorname{Soc}(R)=(0:_R\mathfrak m)\) 满足 \(\dim_k\operatorname{Soc}(R)=1\)。
\end{lemma}
\begin{proof}
证明要对同一个 \(\operatorname{Tor}_c^S(k,R)\) 作两次计算：定义方程的 Koszul 消解把它算成 \(k\)，而环境环正则参数的 Koszul 消解把它识别为基座。

零维 Noether 局部环为 Artin 环，极大理想幂零，所以 \(R=\widehat R\)。写 \(R=S/(f_1,\ldots,f_c)\)，其中 \((S,\mathfrak n)\) 为完备正则局部环，\(f_1,\ldots,f_c\in\mathfrak n\) 为 \(S\) 正则序列。维数公式给出 \(c=\dim S\)。取 \(\mathfrak n\) 的一个极小生成组 \(x_1,\ldots,x_c\)，由\cref{b3:thm:regular-local-characterizations}，它是正则序列。它们在 \(R\) 中的像生成极大理想，即
\[
(x_1,\ldots,x_c)R=\mathfrak nR=\mathfrak m.
\]
且 \(S/\mathfrak n\cong R/\mathfrak m=k\)。

用 \(f\) 的 Koszul 消解计算，\cref{b3:prop:ci-koszul-resolution}给出 \(\operatorname{Tor}_c^S(k,R)\cong k\)。另一方面，用 \(x\) 的 Koszul 复形消解 \(k\)，再与 \(R\) 张量，其最高次同调为
\[
\begin{aligned}
H_c(K(x;S)\otimes_SR)
&=\{a\in R:x_1a=\cdots=x_ca=0\}\\
&=(0:_R\mathfrak m).
\end{aligned}
\]
这里最高项是 \(R\)，微分的各坐标为 \(\pm x_ja\)，且更高一项为零，故最高同调正是所写的共同零化子。两种计算得到同一个 Tor；其一致性由\cref{b3:lem:tor-basic-properties}的平衡性保证。因此基座的维数为一；\(c=0\) 时 \(R=k\)，结论同样成立。
\end{proof}

基座一维只是零维完全交局部环的必要条件，不能反过来作为完全交的充分判据；这个引理也不适用于前面一维的 \(k[\![x,y]\!]/(x^2,xy)\)，该环已由深度小于维数排除完全交性。

\begin{proposition}\label{b3:prop:ci-cm-counterexample-families}
设 \(k\) 为域。
\begin{enumerate}
\item 设 \(n\ge2\) 为整数，\(A=k[\![x]\!]/(x^n)\)。\(A\) 为零维完全交局部环，且交替使用乘 \(x\) 与乘 \(x^{n-1}\) 的复形
\[
\cdots\xrightarrow{x}A\xrightarrow{x^{n-1}}A
\xrightarrow{x}A\longrightarrow k\longrightarrow0
\]
为剩余域 \(k\) 的极小自由 \(A\) 消解。因此
\[
\operatorname{Tor}_i^A(k,k)\cong k\quad(i\ge0),
\qquad\operatorname{pd}_Ak=\infty.
\]
\item 设 \(e\ge2\) 为整数，
\[
B=k[u_1,\ldots,u_e]/(u_1,\ldots,u_e)^2,
\qquad\mathfrak m=(\bar u_1,\ldots,\bar u_e).
\]
则 \(B\) 为零维 CM 局部环，长度为 \(e+1\)，嵌入维数为 \(e\)，且
\[
\operatorname{Soc}(B)=\mathfrak m,
\qquad\dim_k\operatorname{Soc}(B)=e.
\]
因此 \(B\) 不是完全交局部环。
\end{enumerate}
\end{proposition}
\begin{proof}
对第一族，\(k[\![x]\!]\) 是一维正则局部整环，\(x^n\) 为非零非单位，故\cref{b3:prop:given-complete-intersection}说明 \(A\) 为零维完全交局部环。在 \(A\) 中，乘 \(x\) 的核为 \((x^{n-1})\)，乘 \(x^{n-1}\) 的核为 \((x)\)，而 \(A\to k\) 的核也为 \((x)\)，所以所给的列正合。两个乘法系数都在极大理想中，故消解极小；张量 \(k\) 后微分全部为零，得到每个次数的 Tor 都同构于 \(k\)。这些 Tor 在任意高次非零，因此 \(\operatorname{pd}_Ak=\infty\)。当 \(n=2\) 时，所有微分都是乘 \(x\)。

对第二族，\(B\) 的 \(k\) 向量空间基为 \(1,\bar u_1,\ldots,\bar u_e\)。每个元素唯一写成 \(a+v\)，其中 \(a\in k\)、\(v\in\mathfrak m\)；\(a\ne0\) 时由 \(v^2=0\) 得到逆元 \(a^{-1}-a^{-2}v\)，所以 \(\mathfrak m\) 是唯一极大理想。又 \(\mathfrak m^2=0\)，故 \(B\) 为零维局部环，自动为 CM 环。它的所有简单因子均为剩余域 \(k\)，因而长度等于其 \(k\) 维数，即 \(e+1\)；由 \(\mathfrak m/\mathfrak m^2=\mathfrak m\)，嵌入维数为 \(e\)。乘以 \(\bar u_i\) 时有 \(\bar u_i(a+v)=a\bar u_i\)，故同时被全部 \(\bar u_i\) 零化恰好意味着 \(a=0\)。因此基座正是 \(\mathfrak m\)，其维数为 \(e\ge2\)。\cref{b3:lem:artinian-ci-socle}遂排除了所有可能的完全交展示。
\end{proof}

\begin{center}
\begin{minipage}{\textwidth}
\makeatletter
\def\@captype{figure}
\makeatother
\centering
\[
\begin{tikzcd}[column sep=large]
\text{正则} \arrow[r, Rightarrow, "{c=0}"] &
\text{完全交} \arrow[r, Rightarrow, "{\text{正则序列商}}"] &
\text{Cohen--Macaulay}
\end{tikzcd}
\]
\[
\text{Cohen--Macaulay}
\quad\Longleftrightarrow\quad
\text{对每个整数}\;r\ge1\;\text{满足}\;(S_r).
\]
\begin{tabularx}{\textwidth}{@{}>{\centering\arraybackslash}X>{\centering\arraybackslash}X@{}}
完全交而非正则 & CM 而非完全交\\
\(k[\![x]\!]/(x^2)\) & \(k[x,y]/(x,y)^2\)
\end{tabularx}
\caption[正则、完全交与 Cohen--Macaulay 性]{正则、完全交与 Cohen--Macaulay 性。性质关系均限于交换 Noether 局部环；下方两个逆向反例中的 \(k\) 为任意域。}
\label{b3:fig:regular-ci-cm}
\end{minipage}
\end{center}

图中的第二行由\cref{b3:thm:cm-localization}与\cref{b3:def:serre-sk}得到。若 \(R\) 为 CM 局部环，则对每个素理想 \(\mathfrak p\)，局部化后有 \(\operatorname{depth}R_{\mathfrak p}=\dim R_{\mathfrak p}\)，自然满足任意 \((S_r)\)。反之，若 \(R\) 满足所有 \((S_r)\)，则在任一 \(\mathfrak p\) 处选取整数 \(r\ge\max\{1,\dim R_{\mathfrak p}\}\)，其条件便变成
\[
\operatorname{depth}R_{\mathfrak p}\ge\dim R_{\mathfrak p}.
\]
结合深度不超过维数，得到该点的 CM 等式；尤其在极大理想处得到 \(R\) 为 CM 环。这里的“所有 \(r\)”保证在每个点都能取到不小于局部维数的整数，因而恰好要求各局部环的深度达到其维数。

定义只要求存在一个合适的展示，而实际计算通常给出一个指定的正则局部环境。要使两者接上，需要比较不同环境中的关系数。极大理想的第一 Koszul 同调提供了这种比较：多添一个可以消去的环境参数，就同时多出一个零生成元的 Koszul 同调类；扣除这些类以后，剩下的数据只依赖商环。

\begin{theorem}\label{b3:thm:ci-presentation-independence}
设 \((S,\mathfrak n)\) 为交换 Noether 正则局部环，\(I\subsetneq S\) 为理想，\(R=S/I\)。记 \(\mu_S(I)=\dim_{S/\mathfrak n}I/\mathfrak nI\)，即 \(I\) 的最小生成元数。下列条件等价：
\begin{equivalences}
\item \(R\) 是完全交局部环；
\item \(I\) 可由一个 \(S\) 正则序列生成；
\item \(I\) 的任意极小生成组都是 \(S\) 正则序列；
\item \(\mu_S(I)=\operatorname{ht}_S I\)。
\end{equivalences}
\end{theorem}
\symidx{\(\mu_S(I)\)：局部环 \(S\) 中理想 \(I\) 的最小生成元数}
\begin{proof}
记 \(\mathfrak m=\mathfrak n/I\)、\(k=S/\mathfrak n=R/\mathfrak m\)、\(n=\dim S\)、\(e=\operatorname{edim}R\)、\(d=\dim R\)。取 \(\mathfrak m\) 的极小生成组 \(z=(z_1,\ldots,z_e)\)，并记
\[
h(R)=\dim_k H_1(z;R).
\]
\cref{b3:prop:koszul-signs-homotopy}说明 \(\mathfrak m\) 零化这个同调模；它又是有限生成模，因而上式是有限的 \(k\) 维数。这个数不依赖极小生成组：若 \(w\) 是另一组，则 \(w=Uz\) 对某个 \(e\times e\) 矩阵 \(U\) 成立。两组模 \(\mathfrak m^2\) 都是 \(\mathfrak m/\mathfrak m^2\) 的基，所以 \(U\) 模 \(\mathfrak m\) 可逆，\(\det U\) 为单位。\cref{b3:prop:koszul-generator-change}遂给出相应 Koszul 复形的同构。

取 \(\mathfrak n\) 的极小生成组 \(x_1,\ldots,x_n\)。由\cref{b3:thm:regular-local-characterizations}，它们是正则序列，故 \(K(x;S)\) 消解 \(k\)。它们在 \(R\) 中的像生成 \(\mathfrak m\)；从这些像中选出 \(e\) 个，使一次类组成 \(\mathfrak m/\mathfrak m^2\) 的基，重排以后记作 \(z_1,\ldots,z_e\)。Nakayama 引理保证它们生成整个 \(\mathfrak m\)。其余每个像都可写成这些 \(z_i\) 的 \(R\) 线性组合；从相应行减去这个组合，给出 \(R\) 上的可逆矩阵，把整组像变成
\[
(z_1,\ldots,z_e,\underbrace{0,\ldots,0}_{n-e\;\text{项}}).
\]
这里仅在张量后的复形 \(K(x;S)\otimes_SR\) 上换基，不需要把这个矩阵提升为 \(S\) 上的参数变换。由\cref{b3:prop:koszul-generator-change,b3:prop:koszul-tensor}，新增的零生成元贡献微分为零的外代数。总次数一分成 \(H_1(z;R)\otimes_R\bigwedge^0R^{n-e}\) 与 \(H_0(z;R)\otimes_R\bigwedge^1R^{n-e}\)；因 \(H_0(z;R)=R/\mathfrak m=k\)，得到
\[
\operatorname{Tor}_1^S(k,R)
\cong H_1(z;R)\oplus k^{n-e}.
\]
另一方面，对 \(0\to I\to S\to R\to0\) 使用\cref{b3:lem:tor-basic-properties}的 Tor 长正合列。\(S\) 自由，且 \(I\subseteq\mathfrak n\) 使 \(I/\mathfrak nI\to S/\mathfrak n\) 为零，故
\[
\operatorname{Tor}_1^S(k,R)\cong I/\mathfrak nI,
\qquad
\mu_S(I)=h(R)+n-e.
\]
前面高度判据证明中的素理想维数等式还给出
\(\operatorname{ht}_S I=n-d\)。因此，对任意这样的正则局部展示，都有
\[
\mu_S(I)=\operatorname{ht}_S I
\quad\Longleftrightarrow\quad h(R)=e-d.
\]
右边已经不含环境环 \(S\)。

这三个数在完备化下也不变。\cref{b3:thm:completion-dimension}给出 \(\dim\widehat R=d\)；\cref{b3:prop:completion-quotients}把 \(\mathfrak m/\mathfrak m^2\) 识别为 \(\mathfrak m\widehat R/(\mathfrak m\widehat R)^2\)，故嵌入维数仍为 \(e\)。有限自由 Koszul 复形经忠实平坦完备化后，其同调由\cref{b3:thm:completion-flat}变成 \(H_1(z;R)\otimes_R\widehat R\)。这个模被 \(\mathfrak m\) 零化，而完备化的剩余域仍为 \(k\)，所以它与原同调模同构。于是 \(h(\widehat R)=h(R)\)。

若第一项成立，按定义写 \(\widehat R=T/J\)，其中 \(T\) 为完备正则局部环，\(J\) 由长度 \(c\) 的正则序列生成。\cref{b3:prop:given-complete-intersection}给出 \(\operatorname{ht}_TJ=c\)；又有 \(\mu_T(J)\le c\)，而高度定理给出反向不等式，故 \(\mu_T(J)=\operatorname{ht}_TJ\)。把刚得到的关系数等式用于这个展示，得到
\[
h(\widehat R)=\operatorname{edim}\widehat R-\dim\widehat R.
\]
三个数的完备化不变性于是给 \(h(R)=e-d\)，再用于当前展示 \(R=S/I\)，得到第四项。

若第四项成立，任取 \(I\) 的极小生成组，其长度等于 \(\operatorname{ht}_SI\)，\cref{b3:prop:given-complete-intersection}保证这组元素正则，得到第三项。第三项蕴含第二项；第二项又由同一命题蕴含第一项。\(I=0\) 时它由空正则序列生成；\(e=0\) 时 \(z\) 为空，\(R=k\)、\(h(R)=0\)，相同的 Koszul 计算给出 \(\mu_S(I)=n\)，所以这些边界情形也包括在论证内。
\end{proof}

这个比较说明，判定强度取决于是否已经去掉冗余方程。若当前正则局部环境中一个极小生成组不正则，就可以否定环自身的完全交性；仅仅在某个非极小方程组中发现零因子，则还不能作这个结论。

完全交性还应当能从极大理想处传到其他素点。若环本身已经写成正则局部环的正则序列商，直接局部化即可；若只知道完备化有这样的展示，则还要从完备化上的素点下降回原环。这一步需要如下平坦局部下降定理：对交换 Noether 局部环之间的平坦局部同态 \(A\to B\)，若 \(B\) 是完全交局部环，则 \(A\) 也是完全交局部环。这项一般下降定理的证明不由平坦性检测正则序列直接得到；它通过正则局部覆盖及有限 Tor 维数比较两边的定义理想，使用分幂消解。其完整证明见 Stacks Project, Proposition 23.8.4 及 Lemma 23.7.6\footnote{分别见 \url{https://stacks.math.columbia.edu/tag/09Q2} 与 \url{https://stacks.math.columbia.edu/tag/09PX}。这里采用前一命题的下降方向；完整比较还涉及闭纤维的完全交性。本书不在此展开分幂消解与相应有限 Tor 维数判据的证明。}；下面使用这项输入证明局部化结论。

\begin{proposition}\label{b3:prop:ci-localization}
设 \((R,\mathfrak m)\) 为交换 Noether 完全交局部环。对每个 \(\mathfrak p\in\operatorname{Spec}R\)，\(R_{\mathfrak p}\) 都是完全交局部环。因此，完全交局部环也满足\cref{b3:def:local-complete-intersection}中局部完全交环的条件。
\end{proposition}
\begin{proof}
先处理具有实际正则局部展示的情形。若 \(R=S/(f_1,\ldots,f_c)\)，其中 \(S\) 为正则局部环、\(f\) 为 \(S\) 正则序列，取 \(\mathfrak p\) 在 \(S\) 中的原像 \(\mathfrak q\)。由\cref{b3:thm:regular-local-localization}，\(S_{\mathfrak q}\) 正则。局部化的正合性保持每个逐次商上的单射乘法，而全部 \(f_i\in\mathfrak q\) 保证末端商仍非零，因此这个序列在 \(S_{\mathfrak q}\) 上正则。由
\[
R_{\mathfrak p}\cong S_{\mathfrak q}/(f_1,\ldots,f_c)S_{\mathfrak q}
\]
及\cref{b3:prop:given-complete-intersection}，得到 \(R_{\mathfrak p}\) 完全交。这一论证也适用于由\cref{b3:thm:ci-presentation-independence}恢复正则序列核的指定正则局部展示。

一般情形中，定义给出 \(\widehat R=S/(f_1,\ldots,f_c)\)。\cref{b3:thm:completion-flat}保证 \(R\to\widehat R\) 忠实平坦，所以其素谱映射满射。取 \(\mathfrak q\in\operatorname{Spec}\widehat R\)，使 \(\mathfrak q\cap R=\mathfrak p\)。于是
\[
R_{\mathfrak p}\longrightarrow (\widehat R)_{\mathfrak q}
\]
是平坦局部同态：平坦性由局部化保持，目标极大理想的原像为 \(\mathfrak pR_{\mathfrak p}\)。闭纤维非零，\cref{b3:thm:faithfully-flat-criterion}还保证它忠实平坦。前一段已说明 \((\widehat R)_{\mathfrak q}\) 完全交；现在使用所述平坦局部下降定理，得到 \(R_{\mathfrak p}\) 完全交。
\end{proof}

这里没有把 \((R_{\mathfrak p})^{\wedge}\) 与 \((\widehat R)_{\mathfrak q}\) 视为同一个环；两者通过上面的平坦局部比较传递完全交性。一般地，也不能仅由 \(A\to B\) 平坦就把完全交性任意双向传递，向上保持还须对闭纤维加条件。

% !TeX root = ../../Book3.tex
% 纲要标题：### 19.4 超曲面、嵌入分量与重数
% 纲要位置：计划纲要.md，第 2575 行。

\section{超曲面、嵌入分量与重数}
\label{b3:sec:1904}

一个超曲面可以约化或非约化、正规或不正规，却仍然是 Cohen–Macaulay 环。将这些例子与具有嵌入关联素理想的商环比较，可以区分零因子、奇点与深度不足各自造成的现象。


% 纲要内容（仅作写作计划，尚非教材正文）：
%
% * 超曲面与零维局部环
% * 非 CM 的嵌入分量例子；用素理想回避直接证明一维约化 Noether 局部环为 CM 环
% * 参数逐项取幂与参数理想整体取幂的区别；先核算一元幂层及二元二次关系，再用参数个数与全次数的双重归纳证明伴随分次模为多项式模
% * CM 参数理想的精确累计长度公式与重数等于参数商长度；非 CM 情形的对照算例
%
% ---
%

\subsection{超曲面与零维局部环}
\label{b3:subsec:1904:01}

设 \((S,\mathfrak n)\) 为交换 Noether 正则局部环，\(0\ne f\in\mathfrak n\)。商 \(S/(f)\) 称为这里的\term[chaoqumianhuan]{超曲面环}[hypersurface ring]。因为 \(S\) 为整环，\(f\) 是非零因子，所以超曲面环是完全交环，深度和维数都等于 \(\dim S-1\)。是否约化、是否正规、是否正则，则取决于方程的进一步性质。

\ProseParagraphBreak
设 \(k\) 为域。在 \(S=k[\![x,y]\!]\) 中，\(S/(y^2-x^3)\)、\(S/(xy)\) 和 \(S/(x^2)\) 都是一维 CM 环。第一个为整环，给出单个尖点分支；第二个约化，有两个相交分支；第三个含非零幂零元，支集仍只有一个不可约分量。三个环的极大理想模平方都具有 \(k\) 维数二，大于环维数一，故都不正则。CM 性只要求深度达到维数，并不能据此区分约化性、不可约性或这些不同的奇点。

\begin{proposition}\label{b3:prop:zero-dimensional-cm}
每个零维交换 Noether 局部环及其非零有限生成模都是 CM 的。
\end{proposition}
\begin{proof}
设环的极大理想为 \(\mathfrak m\)。它是幂零理想，所以非零有限生成模 \(M\) 有有限长度。取最小的 \(r\ge1\) 使 \(\mathfrak m^rM=0\)，则 \(\mathfrak m^{r-1}M\ne0\) 被 \(\mathfrak m\) 零化。因而不存在属于 \(\mathfrak m\) 的 \(M\) 非零因子，深度为零；模维数也为零。
\end{proof}

零维时，非零有限生成模的基座非零，深度和维数便自动同为零，CM 性没有额外限制。要区分这类环的结构，还须考察长度、基座维数或完全交性；上一节的基座计算表明，零维完全交环仍须满足基座维数为一这一限制。

\subsection{非 Cohen–Macaulay 的嵌入分量}
\label{b3:subsec:1904:02}

设 \(k\) 为域，令
\[
A=k[x,y]_{(x,y)}/(x^2,xy),\qquad\mathfrak m=(\bar x,\bar y).
\]
其根理想为 \((x)\)，故 \(\dim A=1\)。又 \(\bar x\ne0\) 且
\(\mathfrak m\bar x=0\)，所以 \(\operatorname{Ann}_A(\bar x)=\mathfrak m\)。因而 \(\mathfrak m\in\operatorname{Ass}_AA\)，它严格包含极小素理想 \((\bar x)\)，是嵌入关联素理想，并使 \(\operatorname{depth}A=0\)。一维主分量的环 \(A/(\bar x)\cong k[y]_{(y)}\) 本身正则；深度不足来自集中在闭点的子模 \(k\bar x\)。乘法 \(\bar y\) 的核正是这一子模。

\ProseParagraphBreak
\(\bar y\) 是参数，因为 \(A/(\bar y)\cong k[x]/(x^2)\) 长度为二；但它不是非零因子。若换成参数 \(\bar y+a\bar x\)，其中 \(a\in k\)，仍有 \((\bar y+a\bar x)\bar x=0\)。实际上，由 \(\mathfrak m\bar x=0\)，极大理想中的每个元素都是零因子，任何参数都无法成为正则元素。能否选出正则参数取决于环的结构，不能靠更换参数消除这个非零基座元素。

\ProseParagraphBreak
一维约化交换 Noether 局部环 \((R,\mathfrak m)\) 则必为 CM 环。约化性使所有极小分量上同时非零的元素成为非零因子：若一个乘积为零，各整环分量上的消去会迫使另一个因子属于零根。具体地，\(R\) 只有有限个极小素理想，且这些理想都严格包含于 \(\mathfrak m\)。由素理想回避，可取
\[
x\in\mathfrak m\setminus
\bigcup_{\mathfrak p\in\operatorname{Min}R}\mathfrak p.
\]
若 \(xy=0\)，则在每个整环 \(R/\mathfrak p\) 中，\(x\) 的像非零，故 \(y\in\mathfrak p\)。于是 \(y\) 属于所有极小素理想的交，即 \(\sqrt{(0)}\)；约化性使这个交为零。因此 \(x\) 是非零因子，深度至少为一，再由维数上界得到深度等于一。二维约化性本身一般不足以保证还能找到第二个正则元素；\cref{b3:exm:r1-not-s2}将给出两张平面只在闭点相交、深度却只有一的具体环。

\subsection{参数、重数与几何交}
\label{b3:subsec:1904:03}

参数理想给出的有限长度保留了相交时的重数信息。先看交换环 \(R\) 上的模 \(N\)，以及 \(N\) 的非零因子 \(x\in R\)。对每个 \(j\ge0\)，有
\[
N/xN\xrightarrow{\ \cdot x^j\ }x^jN/x^{j+1}N.
\]
它由生成元的形式给出满射；若 \(x^jn=x^{j+1}n'\)，约去非零因子 \(x^j\) 就得 \(n\in xN\)，所以也是单射。当 \(N/xN\) 有有限长度时，\(N/x^aN\) 的幂次滤过有 \(a\) 个这样的层，故长度为 \(a\ell_R(N/xN)\)。多个正则参数的逐项幂次计算，就是在各个参数方向上重复这一过程。

\begin{proposition}\label{b3:prop:cm-parameter-powers-length}
设 \((R,\mathfrak m)\) 为交换 Noether 局部环，\(M\ne0\) 为维数 \(d\) 的\(R\) 上有限生成的 CM 模，\(x_1,\ldots,x_d\in\mathfrak m\) 为 \(M\) 的参数系。对正整数 \(a_1,\ldots,a_d\)，有
\[
\ell_R\bigl(M/(x_1^{a_1},\ldots,x_d^{a_d})M\bigr)
=a_1\cdots a_d\,\ell_R\bigl(M/(x_1,\ldots,x_d)M\bigr).
\]
\end{proposition}
\begin{proof}
幂次参数系仍有同样的根理想，故仍为参数系；\cref{b3:thm:cm-parameters}使它及其任意重排均正则。先把除 \(x_1\) 以外的参数幂商去，得模 \(N\)，其中 \(x_1\) 为非零因子。对每个 \(j\ge0\)，乘 \(x_1^j\) 诱导同构
\[
N/x_1N\longrightarrow x_1^jN/x_1^{j+1}N.
\]
满射来自定义；若 \(x_1^jn=x_1^{j+1}n'\)，可约去非零因子 \(x_1^j\)，故核为零。用这些商滤过 \(N/x_1^{a_1}N\)，长度相加得到因子 \(a_1\)。依次对其余参数做同样的计算。
\end{proof}

在 \(k[\![x,y]\!]\) 中，参数理想 \((x^a,y^b)\) 的商以
\(x^iy^j\ (0\le i<a,\ 0\le j<b)\) 为基，长度为 \(ab\)。它的支集只有原点，但商代数记住了两条带重数坐标轴相交的程度。

\ProseParagraphBreak
当参数个数大于一时，各参数分别取幂通常不同于整个参数理想取幂。例如对域 \(k\)，在 \(k[\![x,y]\!]\) 中有
\[
(x^2,y^2)\subsetneq\mathfrak q^2=(x^2,xy,y^2),
\qquad \mathfrak q=(x,y).
\]
左边把每个参数分别平方，右边包含所有总次数为二的混合项。因此前述长度公式不能直接代入 Hilbert–Samuel 累计函数；后者计算的是 \(M/\mathfrak q^{n+1}M\)，其中 \(\mathfrak q=(x_1,\ldots,x_d)\)。为计算这个商，需要同时控制所有总次数相同的参数单项式。

要证明参数单项式确实形成各层的基，关键是排除新的同次关系。二元二次情形的具体目标是：若 \(x,y\) 为 \(M\) 正则序列，则
\[
ax^2+bxy+cy^2\in(x,y)^3M
\quad\Longrightarrow\quad a,b,c\in(x,y)M
\qquad(a,b,c\in M).
\]
可以直接验证这一点。把右边的三次项写成
\(x^3a'+x^2yb'+xy^2c'+y^3d'\)，并分别从 \(a,b,c\) 中减去
\(xa'+yb'\)、\(yc'\)、\(yd'\)，便得到与原系数模 \((x,y)M\) 相同的系数 \(a_0,b_0,c_0\)，满足
\[
a_0x^2+b_0xy+c_0y^2=0.
\]
模 \(xM\) 后，\(y^2c_0=0\)；\(y\) 在 \(M/xM\) 上为非零因子，所以 \(c_0=xc_1\)。在原等式中约去 \(x\)，得到
\(xa_0+y(b_0+yc_1)=0\)。再模 \(xM\)，得 \(b_0+yc_1=xb_1\)，从而 \(b_0\in(x,y)M\)。代回后再次约去 \(x\)，得 \(a_0=-yb_1\in(x,y)M\)，而 \(c_0=xc_1\) 也属于这个子模。因此原来的三个系数都在 \((x,y)M\) 中。

\ProseParagraphBreak
这个计算先把较高次项吸收到系数中，再用正则性消去同次关系。形式多项式 \(aT_1^2+bT_1T_2+cT_2^2\) 的 \(T_i\) 是形式变量，\(x,y\) 则是作用于 \(M\) 的环元素；只有把 \(T_1,T_2\) 分别代为 \(x,y\) 后，才得到模 \(M\) 中的实际关系。一般证明将对参数个数作外层归纳，再对关系的总次数作内层归纳，把上述消去推广到每一层。

\begin{lemma}[正则序列的伴随分次模]\label{b3:lem:regular-sequence-associated-graded}
设 \((R,\mathfrak m)\) 为交换 Noether 局部环，\(M\ne0\) 为有限生成的 \(R\) 模，\(r\ge0\) 为整数，\(x_1,\ldots,x_r\in\mathfrak m\) 为 \(M\) 正则序列，\(\mathfrak q=(x_1,\ldots,x_r)\)。在 \(\operatorname{gr}_{\mathfrak q}(M)=\bigoplus_{n\ge0}\mathfrak q^nM/\mathfrak q^{n+1}M\) 上，让 \(R/\mathfrak q\) 按原来的标量乘法作用，让形式变量 \(T_i\) 按乘 \(x_i\) 所诱导的次数一映射作用。则规范的分次 \((R/\mathfrak q)[T_1,\ldots,T_r]\) 模映射
\[
\Phi:(M/\mathfrak qM)[T_1,\ldots,T_r]
\longrightarrow\operatorname{gr}_{\mathfrak q}(M),
\qquad
\bar mT^\alpha\longmapsto x^\alpha m+\mathfrak q^{|\alpha|+1}M
\]
为同构。这里 \(\alpha=(\alpha_1,\ldots,\alpha_r)\in\mathbb N^r\)，\(|\alpha|=\sum_i\alpha_i\)，\(T^\alpha=\prod_iT_i^{\alpha_i}\)，\(x^\alpha=\prod_i x_i^{\alpha_i}\)；左边是以 \(M/\mathfrak qM\) 为系数的多项式模，各 \(T_i\) 的次数为一。
\end{lemma}
\begin{proof}
若把 \(m\) 换成同一个 \(M/\mathfrak qM\) 类的另一代表，\(x^\alpha m\) 的差属于 \(\mathfrak q^{|\alpha|+1}M\)，所以 \(\Phi\) 良定义。\(\mathfrak q^nM\) 由所有 \(|\alpha|=n\) 的 \(x^\alpha m\) 生成，故 \(\Phi\) 在每一次数上满射。单射性的证明有两层归纳：外层减少序列长度 \(r\)，内层减少待消去关系的齐次次数 \(n\)。

\(r=0\) 时映射只在次数零为恒等映射，其余次数均为零。\(r=1\) 时，对每个 \(n\ge0\)，映射
\[
M/x_1M\longrightarrow x_1^nM/x_1^{n+1}M,
\qquad \bar m\longmapsto x_1^nm+x_1^{n+1}M
\]
满射。若 \(x_1^nm=x_1^{n+1}m'\)，由 \(x_1\) 在 \(M\) 上为非零因子，可约去 \(x_1^n\)，得到 \(m=x_1m'\)，所以映射也是单射。把各层放在一起便得
\(\operatorname{gr}_{(x_1)}(M)\cong(M/x_1M)[T_1]\)：正则元素的每一幂层都是第零层的一份复制。

现在设 \(r\ge2\)，并假设结论对长度 \(r-1\) 已成立。令 \(J=(x_1,\ldots,x_{r-1})\)。外层归纳给出
\[
\operatorname{gr}_J(M)\cong(M/JM)[T_1,\ldots,T_{r-1}].
\]
右边的每个齐次层是若干份系数模 \(M/JM\) 的直和，\(x_r\) 只作用于这些系数。正则性保证 \(x_r\) 在 \(M/JM\) 上为非零因子，所以它在每层 \(J^sM/J^{s+1}M\) 上的乘法也单射。因此，对 \(s\) 归纳可得
\[
(J^sM:_Mx_r)=J^sM\qquad(s\ge0).
\]
\(s=0\) 时两边都为 \(M\)。若 \(x_rm\in J^sM\) 且 \(s\ge1\)，先由第 \(s-1\) 步得 \(m\in J^{s-1}M\)，再由第 \(s-1\) 层的乘法单射得 \(m\in J^sM\)；反向包含来自 \(J^sM\) 是子模。这个等式说，\(x_rm\) 若已经属于 \(J^sM\)，则 \(m\) 本身就属于 \(J^sM\)，乘 \(x_r\) 不会使元素相对于 \(J\) 的阶数升高。

先消去取值恰为零的关系：若以 \(M\) 为系数的齐次多项式 \(F(T)\) 满足 \(F(x)=0\)，就证明各系数属于 \(\mathfrak qM\)。这里固定 \(r\)，对次数 \(n\) 作内层归纳。\(n=0\) 时唯一系数就是零。若 \(n\ge1\)，写
\[
F(T)=G(T_1,\ldots,T_{r-1})+T_rH(T),
\]
其中 \(G\) 齐次次数为 \(n\)，\(H\) 齐次次数为 \(n-1\)。由
\(x_rH(x)=-G(x_1,\ldots,x_{r-1})\in J^nM\)
及上述冒号等式，得 \(H(x)\in J^nM\)。于是可取以 \(M\) 为系数、齐次次数为 \(n\) 的多项式 \(h(T_1,\ldots,T_{r-1})\)，使
\[
H(x)=h(x_1,\ldots,x_{r-1}).
\]
把 \(h\) 写成 \(\sum_{i<r}T_ih_i\)，每个 \(h_i\) 齐次次数为 \(n-1\)。多项式
\(H-\sum_{i<r}x_ih_i\) 的次数为 \(n-1\)，且在 \(x\) 处取值为零；其中 \(x_i\) 乘在 \(h_i\) 的系数上，因而不改变形式次数。由内层归纳，其系数都在 \(\mathfrak qM\) 中，而 \(\sum_{i<r}x_ih_i\) 的系数本来就在 \(JM\) 中，故 \(H\) 的系数也在 \(\mathfrak qM\) 中。另一方面，\(G+x_rh\) 是前 \(r-1\) 个形式变量的 \(n\) 次齐次多项式，在 \(x_1,\ldots,x_{r-1}\) 处取值为零。由外层归纳的分次同构，它的各系数属于 \(JM\)。减去 \(x_rh\) 的系数，得到 \(G\) 的各系数属于 \(JM+x_rM=\mathfrak qM\)。这就证明了零关系的所有系数都在 \(\mathfrak qM\) 中。

最后，把模 \(\mathfrak q^{n+1}M\) 为零的关系化为刚才的零关系。若 \(F\) 齐次次数为 \(n\)，且 \(F(x)\in\mathfrak q^{n+1}M\)，取齐次次数为 \(n+1\) 的 \(P\)，使 \(P(x)=F(x)\)，并写 \(P=\sum_iT_iP_i\)，各 \(P_i\) 的次数为 \(n\)。于是 \(F-\sum_ix_iP_i\) 齐次次数为 \(n\)，在 \(x\) 处取值为零；\(\sum_ix_iP_i\) 的系数都属于 \(\mathfrak qM\)。由已证结论，\(F\) 的系数也都在 \(\mathfrak qM\) 中。因此 \(\Phi\) 在每一层的核为零，完成外层归纳。
\end{proof}

一般的 Hilbert–Samuel 定理只保证累计长度在充分大的次数后成为多项式；对于 CM 模的参数理想，伴随分次模的上述同构却能从 \(n=0\) 起算出整个累计函数。重数等于参数商长度，是这个精确公式的最高次系数所给出的结果。

\begin{theorem}[CM 模的参数理想重数]\label{b3:thm:cm-parameter-multiplicity}
设 \((R,\mathfrak m)\) 为交换 Noether 局部环，\(M\ne0\) 为维数 \(d\) 的\(R\) 上有限生成的 CM 模，\(x_1,\ldots,x_d\in\mathfrak m\) 为 \(M\) 的参数系，\(\mathfrak q=(x_1,\ldots,x_d)\)。令 \(B=R/\operatorname{Ann}_RM\)，并用 \(\mathfrak q\) 在 \(B\) 中的像计算 Hilbert–Samuel 重数，记为 \(e_{\mathfrak q}(M)\)。则对每个整数 \(n\ge0\)，有
\[
\ell_R(M/\mathfrak q^{n+1}M)
=\binom{n+d}{d}\,\ell_R(M/\mathfrak qM),
\qquad
e_{\mathfrak q}(M)=\ell_R(M/\mathfrak qM).
\]
当 \(d=0\) 时参数系为空，\(\mathfrak q=0\)，两式中的长度与重数均为 \(\ell_R(M)\)。
\end{theorem}
\begin{proof}
参数系保证 \(\mathfrak q\) 的像为 \(B\) 的极大理想准素理想，故\cref{b3:thm:hilbert-samuel}适用。\(R\) 与 \(B\) 给出相同的子模和长度。由\cref{b3:thm:cm-parameters}，参数系为 \(M\) 正则序列，所以上述引理给出
\[
\operatorname{gr}_{\mathfrak q}(M)
\cong(M/\mathfrak qM)[T_1,\ldots,T_d].
\]
若 \(d>0\)，总次数恰为 \(j\) 的参数单项式有 \(\binom{j+d-1}{d-1}\) 个，各对应一份系数模 \(M/\mathfrak qM\)。因此 \(\mathfrak q^jM/\mathfrak q^{j+1}M\) 是这么多份 \(M/\mathfrak qM\) 的直和，从而
\[
\ell_R(\mathfrak q^jM/\mathfrak q^{j+1}M)
=\binom{j+d-1}{d-1}\,\ell_R(M/\mathfrak qM).
\]
沿 \(M/\mathfrak q^{n+1}M\) 的理想幂滤过相加，利用
\(\sum_{j=0}^n\binom{j+d-1}{d-1}=\binom{n+d}{d}\)，得到累计长度公式。其最高次项为
\(\ell_R(M/\mathfrak qM)n^d/d!\)，与 Hilbert–Samuel 重数的归一化比较即得重数公式。若 \(d=0\)，\(M\) 本身具有有限长度，而 \(\mathfrak q=0\)，累计长度恒为 \(\ell_R(M)\)，同样得到结论。
\end{proof}

正则序列给出多项式型伴随分次模的证明可对照\cite[Theorem 16.2, pp. 125--126]{Matsumura1989CommutativeRingTheory}；环的参数理想重数判据见同书\cite[Theorem 17.11, p. 138]{Matsumura1989CommutativeRingTheory}。

一维时，累计长度公式就是 \(\ell_R(M/t^{n+1}M)=(n+1)\ell_R(M/tM)\)。例如对域 \(k\)，在尖点环 \(R=k[\![x,y]\!]/(y^2-x^3)\) 中，\(x\) 为非零因子，\(R/(x)\cong k[y]/(y^2)\)，因而 \(e_{(x)}(R)=2\)。

\ProseParagraphBreak
上一小节的非 CM 环 \(A\) 中，\(A/(\bar y^{n+1})\) 的 \(k\) 基为
\(1,\bar y,\ldots,\bar y^n,\bar x\)，故对每个 \(n\ge0\)，有
\[
\ell_A(A/(\bar y^{n+1}))=n+2=(n+1)+1.
\]
前面的 \(n+1\) 来自一维主分量的基 \(1,\bar y,\ldots,\bar y^n\)；额外的 \(1\) 来自被 \(\bar y\) 零化的有限长度子模 \(k\bar x\)。因此 \(e_{(\bar y)}(A)=1\)，而
\(\ell_A(A/(\bar y))=2\)。重数只读取最高维分量的首项增长，参数商长度还看见这里的零维嵌入部分。这一部分改变了累计长度的常数项，说明非 CM 情形的参数商长度可以严格大于重数；若要把两者相等用作 CM 判据，还须另外陈述和证明其适用条件。


\begin{chaptersummary}
CM 性使深度达到支集维数所允许的最大值。关联素理想维数估计解释了参数系为什么逐项正则；深度局部化不等式、非零因子商的两个降维公式及忠实平坦完备化，使这种性质在常用构造下可控。正则局部环为 CM 环，由正则序列给出的商进一步形成完全交。

\ProseParagraphBreak
这些条件有不同强度。幂零超曲面可以完全交而不正则，基座维数大于一的零维局部环可以 CM 而不完全交；没有嵌入关联素理想也不自动保证高维 CM 性。CM 模的参数理想重数等于参数商的长度，这来自正则序列给出的多项式型伴随分次模；非 CM 例子则显示两者可以不同。后续的正规性判据将从 CM 性取出其中的 \(S_2\) 条件。
\end{chaptersummary}

\BookExercises

\begin{exercise}\label{b3:ex:19-direct-sums}
设 \((R,\mathfrak m)\) 为交换 Noether 局部环，\(M,N\ne0\) 均为\(R\) 上有限生成的 CM 模。证明 \(M\oplus N\) 为 CM 模当且仅当 \(\dim M=\dim N\)。设 \(k\) 为域，在 \(R=k[\![x,y]\!]\) 上给出维数不同的反例。
\end{exercise}
\begin{solution}
有限直和使 \(\operatorname{Ext}_R^i(k,M\oplus N)\cong\operatorname{Ext}_R^i(k,M)\oplus\operatorname{Ext}_R^i(k,N)\)，所以深度等于两者深度的最小值。支集为并，维数等于两者维数的最大值。由于两个模分别 CM，这两个数相等恰好要求维数相同。取 \(M=R\)、\(N=R/(x)\)，维数分别为二与一，直和的深度为一、维数为二。
\end{solution}

\begin{exercise}\label{b3:ex:19-change-annihilator}
设 \(k\) 为域，\(R=k[\![x,y]\!]\)、\(S=R/(x)\)，把 \(M=S\) 分别看成 \(R\) 模和 \(S\) 模。计算这两种底环下的深度与模维数，判断 CM 性及极大 CM 性。说明为什么商去零化模的理想不改变 CM 性，却可能改变极大 CM 性。
\end{exercise}
\begin{solution}
\(S\cong k[\![y]\!]\) 为一维正则局部环，\(M\) 作为 \(S\) 模自由，深度和维数均为一。作为 \(R\) 模，\(y\) 在 \(M\) 上为非零因子，而 \(M/yM\cong k\) 有深度零；故 \(\operatorname{depth}_RM=1\)。其零化理想为 \((x)\)，所以 \(\dim_RM=\dim R/(x)=1\)。两种底环下都为 CM 模，但 \(\dim R=2\)、\(\dim S=1\)，因此 \(M\) 只在 \(S\) 上是极大 CM 模。一般地，商去包含于零化理想的理想不会改变模的深度和维数；极大 CM 性却比较深度与底环维数，后者可以改变。
\end{solution}

\begin{exercise}\label{b3:ex:19-maximal-cm-module}
设 \(k\) 为域，\(R=k[\![x,y,z]\!]\)，令 \(I=(x,y)R\)。构造 \(I\) 的极小自由 \(R\) 消解，计算 \(I\) 与 \(R/(z)\) 的深度和模维数，判断它们是否为 CM 模及极大 CM 模。
\end{exercise}
\begin{solution}
消解为
\[
0\longrightarrow R\xrightarrow{\binom{-y}{x}}R^2
\xrightarrow{(x\quad y)}I\longrightarrow0.
\]
若 \(ax+by=0\)，模 \(xR\) 后利用 \(y\) 在 \(R/(x)\) 上为非零因子，得 \(b=xc\)；再约去 \(x\)，得 \(a=-yc\)。这证明中间正合，左端单射则由整环性得到。微分的条目在极大理想中，故消解极小，\(\operatorname{pd}_RI=1\)。Auslander–Buchsbaum 公式给出 \(\operatorname{depth}_RI=3-1=2\)。非零理想 \(I\) 的零化子为零，所以 \(\dim_RI=3\)，它不是 CM 模。另一方面，\(R/(z)\cong k[\![x,y]\!]\)，深度和模维数均为二，因此为 CM 模，但不是极大 CM 模。
\end{solution}

\begin{exercise}\label{b3:ex:19-local-cm-failure}
设 \(k\) 为域，在 \(R=k[\![x,y,z]\!]\) 中令 \(I=(x,y)R\)。求所有使 \(I_{\mathfrak p}\) 为 \(R_{\mathfrak p}\) 上的 CM 模的素理想 \(\mathfrak p\)。在其余素点计算局部化后的深度和维数，说明非 CM 性可以沿一条闭子集出现，而不只集中在闭点。
\end{exercise}
\begin{solution}
若 \(\mathfrak p\notin V(x,y)\)，则 \(x,y\) 中至少一个在 \(R_{\mathfrak p}\) 中可逆，故 \(I_{\mathfrak p}=R_{\mathfrak p}\)，为正则局部环上的 CM 模。若 \(\mathfrak p\in V(x,y)\)，上一题的消解局部化后仍正合；微分条目 \(x,y\) 都属于 \(\mathfrak pR_{\mathfrak p}\)，故仍为长度一的极小自由消解。记 \(h=\dim R_{\mathfrak p}=\operatorname{ht}\mathfrak p\)。局部化后的非零理想零化子为零，模维数为 \(h\)，而 Auslander–Buchsbaum 公式给出深度 \(h-1\)，所以不是 CM 模。因而 CM 素点恰为 \(D(x)\cup D(y)\)，非 CM 素点恰为 \(V(x,y)\)。由于 \(R/(x,y)\cong k[\![z]\!]\)，后者包含高度二的素理想 \((x,y)\) 和高度三的闭点 \((x,y,z)\)，相应深度与维数分别为 \((1,2)\) 和 \((2,3)\)。
\end{solution}

\begin{exercise}\label{b3:ex:19-quadric-parameters}
设 \(k\) 为域，\(A=k[\![x,y,z]\!]/(xy-z^2)\)。证明 \(x,y\) 的像为参数系，并计算 \(\ell_A(A/(x^a,y^b))\)，其中 \(a,b\ge1\) 为整数。
\end{exercise}
\begin{solution}
\(A\) 是二维 CM 超曲面环，而 \(A/(x,y)\cong k[\![z]\!]/(z^2)\)，长度为二。因此 \(x,y\) 是参数系，自动正则。参数幂长度公式给出所求长度 \(2ab\)。也可用关系 \(z^2=xy\) 把元素唯一写为关于 \(x,y\) 的多项式加 \(z\) 乘这样的多项式，得到 \(x^iy^j,x^iy^jz\)（\(i<a,j<b\)）这 \(2ab\) 个基向量。
\end{solution}

\begin{exercise}\label{b3:ex:19-thickened-line}
设 \(k\) 为域，\(r\ge2\)、\(n\ge1\) 为整数，令 \(A=k[\![x,y]\!]/(x^r)\)。证明 \(A\) 为一维 CM 环，计算 \(\ell_A(A/(y^n))\) 和 \(e_{(y)}(A)\)，并判断它是否正则。
\end{exercise}
\begin{solution}
\(x^r\) 是正则局部整环 \(k[\![x,y]\!]\) 中的非零因子，所以 \(A\) 为一维完全交 CM 环。商 \(A/(y^n)\) 的基为 \(x^iy^j\)，其中 \(0\le i<r,0\le j<n\)，长度为 \(rn\)。因而参数理想 \((y)\) 的重数为 \(r\)。极大理想模平方由 \(x,y\) 给出两个线性独立类，嵌入维数二而环维数一，故不正则。
\end{solution}

\begin{exercise}\label{b3:ex:19-embedded-point-dimension-two}
设 \(k\) 为域，\(A=k[\![x,y,z]\!]/(x^2,xy,xz)\)。求 \(\dim A\)、\(\operatorname{depth}A\) 及 \(\operatorname{Ass}A\)。
\end{exercise}
\begin{solution}
根理想为 \((x)\)，故维数二。\(x\ne0\) 且被整个极大理想零化，故深度零、\(\mathfrak m\in\operatorname{Ass}A\)。有正合列 \(0\to kx\to A\to k[\![y,z]\!]\to0\)，所以关联素理想包含在 \(\{\mathfrak m,(x)\}\) 中；极小素理想 \((x)\) 必为关联素理想，因而恰为这两个。
\end{solution}

\begin{exercise}\label{b3:ex:19-product-cm}
设 \(k\) 为域。判断 \(A=k\times k[t]\) 是否为 CM 环，并求其不同极大理想处的局部维数。说明为什么不能要求这两个维数相同。
\end{exercise}
\begin{solution}
每个素理想来自一个因子。来自 \(k\) 的唯一素点处局部环为域，维数零；来自 \(k[t]\) 的极大理想处局部环为 DVR，维数一。它们都是 CM 局部环，由\cref{b3:thm:cm-localization}，检验所有极大理想便足够，所以 \(A\) 为 CM 环。\(k[t]\) 的一般点 \((0)\) 处进一步局部化后为域，也直接符合定义。非局部 CM 性逐点比较深度和维数，并不比较不同连通分支的维数。
\end{solution}

\begin{exercise}\label{b3:ex:19-completion-cusp}
设 \(k\) 为域，令 \(R=k[x,y]_{(x,y)}/(y^2-x^3)\)。写出 \(\widehat R\)，并在原环与完备环中分别求参数 \(x\) 所给商的长度。由此判断 CM 性是否改变。
\end{exercise}
\begin{solution}
有限商与完备化相容，故 \(\widehat R=k[\![x,y]\!]/(y^2-x^3)\)。原环是二维正则局部整环 \(k[x,y]_{(x,y)}\) 模非零因子 \(y^2-x^3\) 所得的商，完备环则是二维正则局部整环 \(k[\![x,y]\!]\) 模同一非零因子所得的商，所以两者都是一维 CM 环。它们商去 \(x\) 后均为 \(k[y]/(y^2)\)，长度二；也可由\cref{b3:thm:cm-completion}比较 CM 性。完备环并不等于原环，这里的对应来自有限阶商保持。
\end{solution}

\begin{exercise}\label{b3:ex:19-two-equation-koszul}
设 \(k\) 为域，在 \(S=k[\![x,y,z]\!]\) 中令 \(A=S/(xy,z^2)\)。证明给定两个方程形成正则序列，写出 \(A\) 的极小自由 \(S\) 消解，并计算 \(A/(x+y)\) 的长度。
\end{exercise}
\begin{solution}
\(xy\) 是 \(S\) 非零因子。商 \(S/(xy)\cong\bigl(k[\![x,y]\!]/(xy)\bigr)[\![z]\!]\) 中乘 \(z^2\) 单射，因为它只把级数系数向后移两位，所以序列正则。消解为
\[
0\longrightarrow S\xrightarrow{\binom{-z^2}{xy}}S^2
\xrightarrow{(xy\quad z^2)}S\longrightarrow A\longrightarrow0.
\]
微分系数都在极大理想中，所以极小。商去 \(x+y\) 后得到 \(k[\![x,z]\!]/(x^2,z^2)\)，基为 \(1,x,z,xz\)，长度四。这也说明 \(x+y\) 是一维 CM 环 \(A\) 的正则参数。
\end{solution}

\begin{exercise}\label{b3:ex:19-redundant-equations}
设 \(k\) 为域。在 \(S=k[\![x,y,z]\!]\) 中，理想 \(I=(xy,xz)\) 由两个元素生成。求 \(\dim S/I\)，指出给定序列为何不正则，并写出一条两生成元间的非平凡关系。
\end{exercise}
\begin{solution}
\(I=(x)\cap(y,z)\)，所以商的两个不可约分量维数分别为二和一，整体维数二。若两个元素正则，维数应从三降到一，已经矛盾。直接看，在 \(S/(xy)\) 中，非零的 \(y\) 被 \(xz\) 零化，故第二项不正则。关系为 \(z(xy)-y(xz)=0\)。
\end{solution}

\begin{exercise}\label{b3:ex:19-square-zero-not-ci}
设 \(k\) 为域，令 \(B=k[x,y]/(x^2,y^3)\)、\(A=B/(xy)\)。求两环的长度、深度、嵌入维数及基座，判断完全交性。由此说明完全交环的任意商未必仍是完全交环。
\end{exercise}
\begin{solution}
\(B\) 的基为 \(1,x,y,xy,y^2,xy^2\)，长度六；\(A\) 的基为 \(1,x,y,y^2\)，长度四。它们均为零维局部环，深度零，极大理想模平方都以 \(x,y\) 的像为基，嵌入维数为二。在 \(B\) 中同时被 \(x,y\) 零化的元素恰为 \(k\cdot xy^2\)，而在 \(A\) 中恰为 \(kx\oplus ky^2\)，故基座维数分别为一和二。\(x^2,y^3\) 在正则局部环 \(k[x,y]_{(x,y)}\) 中形成正则序列：商去 \(x^2\) 后，\(y\) 仍为非零因子。因此 \(B\) 为完全交环；\(A\) 的二维基座则由\cref{b3:lem:artinian-ci-socle}排除了完全交性。此处再商去的 \(xy\) 是零因子，不能应用正则序列商的保持性质。
\end{solution}

\begin{exercise}\label{b3:ex:19-periodic-hypersurface}
设 \(k\) 为域，\(n\ge2\)、\(1\le s<n\) 为整数，\(A=k[\![x]\!]/(x^n)\)、\(M=A/(x^s)\)。构造 \(M\) 的极小自由 \(A\) 消解，并计算各阶 \(\operatorname{Tor}_i^A(M,M)\) 的 \(k\) 维数。
\end{exercise}
\begin{solution}
交替使用乘 \(x^s\) 与乘 \(x^{n-s}\)：
\[
\cdots\xrightarrow{x^s}A\xrightarrow{x^{n-s}}A
\xrightarrow{x^s}A\longrightarrow M\longrightarrow0.
\]
两种乘法的核分别为 \((x^{n-s})\)、\((x^s)\)，所以复形正合；两个指数都正，微分在极大理想中，故消解极小。与 \(M\) 张量后，奇数次微分 \(x^s\) 为零，偶数次微分为乘 \(x^{n-s}\)。因此
\[
\operatorname{Tor}_0^A(M,M)\cong M,\qquad
\operatorname{Tor}_{2j+1}^A(M,M)\cong M/x^{n-s}M\quad(j\ge0),
\]
\[
\operatorname{Tor}_{2j}^A(M,M)\cong(0:_Mx^{n-s})\quad(j\ge1).
\]
在基 \(1,x,\ldots,x^{s-1}\) 上，乘 \(x^{n-s}\) 的秩为 \(\max\{0,2s-n\}\)，故零次 Tor 的维数为 \(s\)，所有正次 Tor 的维数均为 \(\min\{s,n-s\}\)。特别地，\(M\) 的投射维数无穷。
\end{solution}

\begin{exercise}\label{b3:ex:19-pd-codimension}
设 \(k\) 为域，\(S=k[\![x,y,z]\!]\)，令 \(M=S/(x^2,xy)\)、\(N=S/(x^2,yz)\)。构造两者的极小自由 \(S\) 消解，计算投射维数、深度和模维数。说明相同的环境环消解长度为何可能对应不同的 CM 性。
\end{exercise}
\begin{solution}
\(M\) 的消解为
\[
0\longrightarrow S\xrightarrow{\binom{-y}{x}}S^2
\xrightarrow{(x^2\quad xy)}S\longrightarrow M\longrightarrow0.
\]
若 \(ax^2+bxy=0\)，整环性先约去 \(x\)，得 \(ax+by=0\)。模 \(xS\) 后，\(y\) 为非零因子，故 \(b=cx\)；再在原等式中约去 \(x\)，得 \(a=-cy\)。这证明中间正合，左端单射由整环性得到。\(N\) 则由正则序列 \(x^2,yz\) 的 Koszul 复形消解：商去 \(x^2\) 后的环在 \(k[\![y,z]\!]\) 上以 \(1,x\) 为自由基，故乘 \(yz\) 单射。两条消解的非零自由项秩都为 \(1,2,1\)，微分条目都在极大理想中，因而极小，投射维数都为二。Auslander–Buchsbaum 公式给出两模的深度均为一。

\(\sqrt{(x^2,xy)}=(x)\)，所以 \(\dim_SM=2\)；\((x^2,yz)\) 的极小素理想为 \((x,y)\)、\((x,z)\)，所以 \(\dim_SN=1\)。因此 \(M\) 非 CM，而 \(N\) 为 CM。前者的消解长度二超过支集余维一，后者的消解长度恰等于支集余维二；判断 CM 性须比较这两个数，而不能单独看消解长度。
\end{solution}

\begin{exercise}\label{b3:ex:19-plane-crossing-parameter}
设 \(k\) 为域，\(a,b\ge1\) 为整数，令 \(A=k[\![x,y,z]\!]/(xy)\)。证明 \(x+y,z\) 为参数系，并求
\(\ell_A(A/((x+y)^a,z^b))\)。
\end{exercise}
\begin{solution}
该环为二维 CM 超曲面环。商去 \(x+y,z\) 得 \(k[\![x]\!]/(x^2)\)，长度二，所以这两个元素为参数系并且正则。参数幂长度公式给出 \(2ab\)。即使环有两个不可约分量，参数仍可选为非零因子；单独选 \(x\) 则被非零的 \(y\) 零化。
\end{solution}

\begin{exercise}\label{b3:ex:19-semigroup-parameter-free}
设 \(k\) 为域，令 \(A=k[\![t^2,t^5]\!]\subseteq k[\![t]\!]\)。把它写成 \(k[\![t^2]\!]\) 上的有限秩自由模，求 \(\ell_A(A/(t^2))\)，并判断 \(A\) 的 CM 性与正则性。
\end{exercise}
\begin{solution}
半群 \(\langle2,5\rangle\) 的元素是全部非负偶数与至少五的奇数。故每个级数唯一写成 \(f(t^2)+t^5\cdot g(t^2)\)，于是 \(A=k[\![t^2]\!]\oplus t^5k[\![t^2]\!]\)，为秩二自由模。商去 \(t^2\) 后基为 \(1,t^5\)，长度二。\(A\) 为一维整环，\(t^2\) 非零因子，故 CM。极大理想由 \(t^2,t^5\) 最小生成，平方中没有阶二或阶五的项，所以嵌入维数二而非正则。
\end{solution}
